\section{Convex compactness of terminal gains}
We henceforth identify the general-graph terminal sets with the gain-model sets using \contractref{M1}--\contractref{M3}.
Clause \contractref{M4} supplies a.e. measurability and the admissible lower bound at the terminal value. Linear operations belong to the model structure \contractref{M0}.

\subsection{NFLVR and boundedness in probability}
A set $D\subseteq L^0$ is bounded in probability if for every $\varepsilon>0$
there exists $R>0$ with
\[
 \sup_{f\in D}\Pp(|f|>R)\leq\varepsilon.
\]
\begin{proposition}[Delbaen--Schachermayer, Proposition 3.1]\label{prop:k1-bip}
Under NFLVR, $K_1$ is bounded in probability.
\end{proposition}
\begin{proof}
Otherwise, since its elements are at least $-1$, choose $1\leq a_n\uparrow\infty$
and $G_n\in K_1$ with $\Pp(G_n>a_n)>\varepsilon$ for one $\varepsilon>0$.
Put
\[
 f_n=\min(G_n/a_n,1).
\]
Positive scaling and subtraction of a nonnegative claim give $f_n\in C_0$.
Moreover,
\[
 -a_n^{-1}\leq f_n\leq1,\qquad
 \E f_n\geq\Pp(G_n>a_n)-a_n^{-1}>\varepsilon-a_n^{-1}.
\]
Bounded Komlós convexification gives forward convex averages $W_nf$
converging almost surely to some $0\leq f\leq1$. Since $a_n$ increases,
these averages still lie between $-a_n^{-1}$ and $1$, and their expectations
are at least $\varepsilon-a_n^{-1}$. Dominated convergence yields
$\E f\geq\varepsilon$, hence $\Pp(f>0)>0$.
Choose $b>0$ with $\Pp(f\geq b)>0$. Egorov's theorem provides a measurable
subset $A\subseteq\{f\geq b\}$ of positive probability on which $W_nf\to f$
uniformly. Define
\[
 U_n=1_AW_nf-a_n^{-1}1_{A^c},\qquad U=1_Af.
\]
We have $U_n\leq W_nf$, so solidity of $C_0$ and boundedness give $U_n\in C$.
Also $0\ne U\in L^\infty_+$ and
\[
 \|U_n-U\|_\infty
 \leq\max\left\{\sup_{\omega\in A}|(W_nf)(\omega)-f(\omega)|,a_n^{-1}\right\}
 \longrightarrow0.
\]
This contradicts NFLVR.
\end{proof}

\subsection{Forward convex limits}
A forward convex combination of $(x_n)$ is
\[
 y_n=\sum_{k=n}^{N_n}\alpha_{n,k}x_k,\qquad
 \alpha_{n,k}\geq0,\qquad\sum_{k=n}^{N_n}\alpha_{n,k}=1.
\]
The next proposition corresponds to the case of \cite[Lemma A1.1]{DS} in which boundedness in probability of the convex hull ensures a finite limit.

\begin{proposition}[Komlós compactness with a lower bound]\label{prop:forward}
Let $D\subseteq L^0$ be convex and bounded in probability. A sequence $G_n\in D$
with a common constant lower bound has forward convex combinations converging
almost surely to a finite random variable.
\end{proposition}
\begin{proof}
Translate by a constant to make the sequence nonnegative. On each tail convex
hull choose approximate maximizers $Y_n$ of $\E[1-e^{-Y}]$.
The midpoint identity for the exponential implies that $e^{-Y_n/2}$ is
$L^2$-Cauchy. Extract indices $n_j\geq j$ such that
$e^{-Y_{n_j}/2}\to q$ almost surely. If $q=0$ on a set of positive probability,
$Y_{n_j}$ tends to infinity there, contradicting boundedness in probability
of the convex set. Thus $q>0$ almost surely. Set $\widetilde G_j=Y_{n_j}$.
Then $\widetilde G_j\to-2\log q$ almost surely, and its convex weights
are supported on $k\geq n_j\geq j$, as required for a forward combination
of the original sequence. Finally undo the initial constant translation.
\end{proof}
Boundedness in probability of a sequence alone does not imply boundedness of
its growing convex hulls. We use convexity of $K_1$ and boundedness of the
entire set $K_1$.

\subsection{Maximal terminal claims}
The next proposition supplies the maximal element of \cite[Lemma 4.3]{DS}. Here convex compactness and maximization of a functional replace transfinite induction.

\begin{proposition}\label{prop:maximal}
If $D\subseteq L^0$ is convex, bounded in probability, and sequentially closed
in probability, each $g\in D$ is dominated by an almost-sure maximal $h\in D$.
\end{proposition}
\begin{proof}
Maximize $\E[1-e^{-(f-g)}]$ over $D_g=\{f\in D:g\leq f\text{ a.s.}\}$.
The common lower bound $g$ need not be constant. Apply
Proposition~\ref{prop:forward} to the nonnegative convex set $D_g-g$, which is
bounded in probability, and to an approximate maximizing sequence.
Add $g$ back to the convex averages. Their almost-sure limit $h$ belongs to
$D_g$ by sequential closedness. Concavity preserves the tail lower bounds on
the objective, and bounded convergence shows that $h$ attains its supremum.
If $k\in D$, $k\geq h$, and $\Pp(k>h)>0$, strict monotonicity of $1-e^{-x}$
contradicts this maximizing property.
\end{proof}
The probability closure $\clprob{K_1}$ has these three properties.
In particular, any forward convex limit $g$ of claims from $K_1$ is dominated
by a maximal $h\in\clprob{K_1}$.
